\documentclass[12pt]{extarticle}
\def\activityid{F04-X-v3}\usepackage[letterpaper,margin=.65in,top=.60in,bottom=.65in]{geometry}
\usepackage{fix-cm}
\usepackage[T1]{fontenc}\usepackage[default]{sourcesanspro}
\usepackage{microtype,tikz,array,tabularx,fancyhdr,amsmath,amssymb}
\usepackage[hidelinks]{hyperref}\usetikzlibrary{calc,arrows.meta}
\setlength{\parindent}{0pt}\setlength{\parskip}{7pt}\setlength{\headheight}{0pt}\setlength{\footskip}{26pt}
\setlength{\emergencystretch}{2em}\pagestyle{fancy}\fancyhf{}\renewcommand{\headrulewidth}{0pt}\renewcommand{\footrulewidth}{.4pt}
\fancyfoot[L]{\scriptsize Bellingham Math Circle / Week 4 / \activityid}\fancyfoot[R]{\scriptsize \thepage}
\definecolor{ink}{gray}{.12}\definecolor{soft}{gray}{.95}\color{ink}
\newcommand{\PageTitle}[2]{{\fontsize{10}{12}\selectfont\bfseries WEEK 4\quad ROUND AND ROUND\hfill #1\par}\vspace{3pt}{\fontsize{26}{29}\selectfont\bfseries #2\par}\vspace{2pt}}
\newcommand{\NameLine}{{\small Name: \rule{2.65in}{.4pt}\hfill Date: \rule{1.35in}{.4pt}\par}\vspace{4pt}}
\newcommand{\Task}[2]{\par\vspace{3pt}{\large\bfseries #1}\par #2\par}
\newcommand{\Subhead}[1]{\par\vspace{3pt}{\large\bfseries #1}\par}
\newcommand{\RuleBox}[1]{\par\begingroup\setlength{\fboxsep}{8pt}\colorbox{soft}{\parbox{\dimexpr\linewidth-16pt\relax}{#1}}\endgroup\par}
\newcommand{\WriteBox}[2]{\par\textbf{#1}\par\vspace{2pt}\begin{tikzpicture}\draw[black!40,line width=.5pt] (0,0) rectangle (\dimexpr\linewidth-.5pt\relax,#2);\end{tikzpicture}\par}
\newcommand{\StudentFooter}{\vfill{\small Try it with objects. Explain by pointing, talking, or drawing.}\par}
\newcommand{\LampRing}[3]{\begin{tikzpicture}[x=1in,y=1in]
\foreach \i in {1,...,#1}{\coordinate (v\i) at ({90-360*(\i-1)/#1}:#2);}
\foreach \i in {1,...,#1}{\pgfmathtruncatemacro{\j}{mod(\i,#1)+1}\draw[thick] (v\i)--node[midway,fill=white,inner sep=2pt,font=\small]{\i\j}(v\j);}
\foreach \i in {1,...,#1}{\node[circle,draw,fill=white,minimum size=.52in,inner sep=0pt,font=\bfseries] at(v\i){\i};}
\foreach \i in {#3}{\node[circle,draw,fill=ink,text=white,minimum size=.52in,inner sep=0pt,font=\bfseries] at(v\i){\i};}
\end{tikzpicture}}
\newcommand{\Lights}[1]{\begin{tikzpicture}[x=.55in,y=.55in,baseline=-.10in]\foreach \v [count=\i] in {#1}{\ifnum\v=1\fill (\i,0) circle(.16in);\else\draw (\i,0) circle(.16in);\fi}\end{tikzpicture}}
\newcommand{\ClockRing}[2]{\begin{tikzpicture}[x=1in,y=1in]\draw[black!30] (0,0) circle(#2);\pgfmathtruncatemacro{\n}{#1-1}\foreach\i in {0,...,\n}{\node[circle,draw,fill=white,minimum size=.35in,inner sep=0pt] at({90-360*\i/#1}:#2){\i};}\draw[-{Stealth},thick] (-.35,.15) arc(155:25:.38);\node[font=\small] at(0,-.18){clockwise};\end{tikzpicture}}
\newcommand{\Key}[2]{\begin{center}\renewcommand{\arraystretch}{1.55}\begin{tabular}{c|*{#1}{c}}in&#2\end{tabular}\end{center}}


% v3 student pages use one running header and numbered tasks only.
\newcommand{\StudentHeader}[1]{{\fontsize{10}{12}\selectfont\bfseries Week 4 / Ring shifts / #1\par}\vspace{10pt}}
\newcommand{\Problem}[2]{\par\vspace{4pt}\textbf{Problem #1:} #2\par}
\newcommand{\Space}[1]{\begin{tikzpicture}\draw[black!35] (0,0) rectangle (\dimexpr\linewidth-.5pt\relax,#1);\end{tikzpicture}\par}
\newcommand{\Record}[2]{#1\par\Space{#2}}

\begin{document}
\StudentHeader{Extra / Grades 6--7}

\Problem{1}{Lay out cards 0--7, leftmost first. Cut into two equal rows without reversing either half. Interleave one card at a time, starting with the first half. The first shuffle is shown. Rebuild it with cards, then repeat the same rule until the original row first returns. Record every row.}
\begin{center}\renewcommand{\arraystretch}{2}\begin{tabular}{c|p{4.6in}}Shuffle&Card labels, left to right\\\hline0&0 1 2 3 4 5 6 7\\1&0 4 1 5 2 6 3 7\\2&\\3&\\4&\\\end{tabular}\end{center}
\Problem{2}{Number places 0--7 from left to right on a strip of paper. Put the original row above it. Track card 1 through the shuffles, recording its place each time. Reset and track card 3, then card 7.}
\Space{1.7in}

\newpage
\StudentHeader{Extra / Grades 6--7}

\Problem{3}{Reset to 0--7. Shuffle once. Fill where each old position moves. A card's label stays the same; its position can change. For example, the card at old position 4 moves to position 1.}
\begin{center}\renewcommand{\arraystretch}{2}\begin{tabular}{c|*{8}{p{.42in}}}Old position&0&1&2&3&4&5&6&7\\\hline New position&&&&&1&&&\end{tabular}\end{center}
\Problem{4}{Draw arrow loops for this position map, starting at 0 and then at each unused position. Use the loops to predict the first return of the whole row. Check against page 1.}
\Space{1.25in}
\Problem{5}{For positions 0--6, double the old position and subtract 7 if needed to leave a number from 0 to 6. Compare with your measured new positions. Treat position 7 separately.}
\begin{center}\renewcommand{\arraystretch}{1.8}\begin{tabular}{c|*{7}{p{.45in}}}Old position&0&1&2&3&4&5&6\\\hline Double&&&&&&&\\Remainder after 7&&&&&&&\end{tabular}\end{center}
\Problem{6}{Starting with 1, repeatedly double and keep the remainder after division by 7. Record until 1 returns. Compare with the route of old position 1.}
\Space{1in}

\newpage
\StudentHeader{Extra / Grades 6--7}

\Problem{7}{Use six cards 0--5. Repeat the same out-shuffle until the full row returns. Record rows. Then test whether the position rule is doubling with remainder after division by 5, with position 5 fixed.}
\Space{1.75in}
\Problem{8}{For sixteen cards 0--15, predict the return by starting with 1, repeatedly doubling, and taking the remainder after division by 15. Check with cards or draw the position loops; keep position 15 fixed separately.}
\Space{1.75in}
\Problem{9}{Compare the first return times for 6,8,16 cards. Record one comparison showing whether a larger deck must take more shuffles to return.}
\Space{1in}

\end{document}
